% answers_short.tex —— 1985 年·上海卷·文史类 · 简短答案速查
% 由 tools/make_short_answers.py 自 answers.tex 的 \gksolhead/\gkanswer 提取，
% 并与 1985/1985_数学真题_上海文_答案要点.md 逐题比对后生成，勿手改（改源 MD 后重跑）。

\begin{gknotice}
\textbf{使用说明}：本节只给最终答案便于快速核对，完整推导见解析版。答案为编写者独立推导，\textbf{非官方标准答案}。
\end{gknotice}

\gkpart{大题}
\begin{itemize}[leftmargin=2.6em]
  \item[\textbf{1.}] (1) $(-3,-1)$；(2) $y=(x+1)^8\ (x\ge-1)$；(3) $\dfrac\pi2$；(4) $29$；(5) $x=1$；(6) $\left(\dfrac12\right)^{1/7}<1$；(7) $(x-2)^2+(y-4)^2=5$；(8) $\dfrac{3\pi}{4}$；(9) $\dfrac{2(1-a)}a$；(10) $\pi$
  \item[\textbf{2.}] (1) （C）；(2) （D）；(3) （D）；(4) （B）；(5) （A）
  \item[\textbf{3.}] (1) 直线 $y=\dfrac23x+2$；(2) 双曲线 $\dfrac{y^2}{4}-\dfrac{x^2}{9}=1$
  \item[\textbf{4.}] (1) $0$；(2) 15 种取法，60 个四位偶数；(3) $\dfrac{16\pi}{3}$
  \item[\textbf{5.}] $PB\perp b$。
  \item[\textbf{6.}] (1) 见解析正文；(2) $\dfrac{\sqrt3-\sqrt2}{2\sqrt3}$
  \item[\textbf{7.}] 见解析正文。
  \item[\textbf{8.}] (1) $x_1=\dfrac{4x_0+2y_0}{5}$，$y_1=\dfrac{2x_0+y_0}{5}$；(2) $x=\dfrac{3x_0+2y_0}{4}$，$y=\dfrac{x_0}{2}$；(3) $x^2-y^2+2x-2y+8=0$，即 $(y+1)^2-(x+1)^2=8$，是一条等轴双曲线
\end{itemize}
